% =============================================================================
% chapter1.tex  --  Chapter 1: Introduction
% Arnold K. Seong, UCI Ph.D. Dissertation, 2026
% =============================================================================

\chapter{Introduction}
\label{chap:intro}

Variable selection is the task of identifying which
predictors in a candidate set are genuinely relevant to an outcome of
interest and, just as importantly, screening out those which, conditional on the selected predictors, are either not associated with the outcome or whose magnitude of impact is below a determined threshold for scientific relevance.  This task is increasingly important in the age of high-dimensional data regimes: our ability to collect data about many different factors makes it ever more difficult to identify which of these factors are truly important to the outcome that we care about.  Whether our goal is understanding which genetic variants are associated with a disease \citep{tibshirani1996lasso}, which economic or demographic indicators predict future wages, or which neural signals encode a cognitive process, the ability to separate signal from noise is imperative to scientific understanding.

At the same time, we also require more flexible modeling tools to describe complex relationships beyond linear associations.  Modeling neuronal cluster activation in response to timed stimuli, for example, or identifying not only if there are differences in wages between demographic groups \textit{but also when} those wage gaps are present, requires models that have the capacity to learn and express such relationships.  This is, however, a double-edged sword: more flexible models increase the risk of overfitting and finding specious relationships between predictor and outcome, again necessitating principled variable selection methods.  
This dissertation develops Bayesian methods for variable selection in \emph{semi-parametric} settings, where the relationship between predictors and the outcome is only partially specified and may include nonlinear functional components. The two main contributions---a procedure
that tests for local effects in functional data via variable selection and modified basis projections (Chapter~\ref{chap:semipar}) and a horseshoe-prior-augmented Bayesian neural network for input variable selection (Chapter~\ref{chap:horseshoe}) --- are united by a common concern: 
achieving principled, interpretable, and uncertainty-aware variable selection when the model's functional form is flexible or unknown.

The remainder of this chapter provides background needed to situate
these contributions. 
Section~\ref{sec:intro-vs} motivates variable selection and reviews its principal challenges.
Section~\ref{sec:intro-freq} surveys modern frequentist approaches, while
Section~\ref{sec:intro-bayes} introduces Bayesian variable and model selection, contrasting its goals, mechanics, and advantages with the frequentist paradigm. Section~\ref{sec:intro-semipar} extends these ideas to semi-parametric and nonparametric settings.
Section~\ref{sec:intro-data} describes the two benchmark datasets used in
Chapter~\ref{chap:semipar}. 
Section~\ref{sec:intro-outline} provides a
detailed outline of the dissertation.

% -----------------------------------------------------------------------------
\section{The Variable Selection Problem}
\label{sec:intro-vs}
% -----------------------------------------------------------------------------

Let $y \in \R$ denote a scalar response and let
$\bx = (x_1, \ldots, x_p)^{\top} \in \R^p$ denote a vector of $p$ candidate
predictors observed on $n$ units. In the classical linear regression setting, the
data-generating model is
\begin{equation}
  y_i = \beta_0 + \bx_i^{\top}\bbeta^* + \epsilon_i,
  \qquad \epsilon_i \overset{\mathrm{iid}}{\sim} N(0, \sigma^2),
  \quad i = 1,\ldots,n,
  \label{eq:linear-model}
\end{equation}
where $\bbeta^* \in \R^p$ is the true coefficient vector. Variable
selection amounts to recovering the \emph{true active set}
$\mathcal{S}^* = \{j : \beta_j^* \neq 0\}$, i.e. the non-zero coefficients signifying a linear relationship between covariate $x_i$ and the response $y$

The problem is trivially solved under the linearity assumption when $p$ is small relative to $n$: ordinary least squares (OLS) is consistent, and classical methods identify non-zero coefficients with controlled type-I error. When $p$ is large, however, three compounding difficulties arise.

\begin{enumerate}
  \item \textbf{Dimensionality.} When $p \geq n$, OLS is not even identifiable; the system of equations is under-determined, leading to infinite possible coefficient estimates or requiring penalization to obtain usable results.

  \item \textbf{Multiplicity.} Testing $p$ hypotheses simultaneously
  inflates the family-wise error rate unless explicit multiplicity
  corrections are applied, which in turn reduce power.

  \item \textbf{Model uncertainty.} With $2^p$ possible subsets of
  predictors, exhaustive evaluation quickly becomes computationally
  infeasible.  Forward/backward stepwise heuristics do not control for multiplicity, nor do they search the model space in a principled way.
\end{enumerate}

Variable selection becomes considerably more
involved when the mean function is nonlinear or only partially specified.
In a \emph{generalized additive model} (GAM) \citep{hastie1990gam}, for
example, the response depends on predictors through a sum of smooth
functions:
\begin{equation}
  g\bigl(\E[y_i \mid \bx_i]\bigr)
  \;=\; f_0 \;+\; \sum_{j=1}^{p} f_j(x_{ij}),
  \label{eq:gam}
\end{equation}
where $g$ is a known link function and each $f_j$ is an unknown smooth
function. Here, variable selection means determining which $f_j$ are
identically zero---a functional hypothesis rather than a scalar one, requiring tests over groups of coefficients.  Moreover, one may wish to determine whether a non-zero $f_j$ exerts its effect only over a \emph{localized} region of the predictor's support.  Answering this question dramatically increases the number of required hypothesis tests, and, as the work presented in Chapter~\ref{chap:semipar} demonstrates, models employing standard basis expansions such as splines do not yield asymptotically consistent estimators for local effects.

Last, in deep neural networks, understanding which input
features $x_j$ are truly related to the response becomes even more difficult, as the functional form of the relationship is implicitly determined by a myriad of learned parameters whose relationships are difficult to disentangle --- embedded within affine functions and connected by non-linear activations --- even in small network configurations.  

% -----------------------------------------------------------------------------
\section{Frequentist Approaches to Variable Selection}
\label{sec:intro-freq}
% -----------------------------------------------------------------------------

Classical frequentist variable selection falls into two broad families:
subset-selection procedures and penalized-likelihood methods.

\subsection{Subset Selection and Stepwise Procedures}
\label{subsec:subset}

Best-subset selection enumerates all $2^p$ candidate models, fits each by
OLS, and selects the model that minimizes an information criterion such as
AIC \citep{akaike1974aic} or BIC \citep{schwarz1978bic}. When $p$ is
small, this finds the global minimum of the criterion. But the
computational cost grows exponentially with $p$, making exhaustive search
impractical even for moderate $p$.

Forward, backward, and stepwise greedy procedures reduce the search to
$O(p^2)$ model evaluations but sacrifice optimality guarantees. These
algorithms are well known to produce unstable selections: small
perturbations to the data can change the selected model drastically
\citep{breiman1996stability}. Moreover, reported $p$-values following a
stepwise search are not valid in the classical sense, because the same
data were used both to select and to test the model \citep{berk2013valid}.

\subsection{Penalized Likelihood: LASSO and Its Extensions}
\label{subsec:lasso}

A landmark development in variable selection was the \emph{least absolute
shrinkage and selection operator} (LASSO) of \citet{tibshirani1996lasso},
which solves
\begin{equation}
  \hat{\bbeta}^{\mathrm{LASSO}}
  \;=\;
  \arg\min_{\bbeta \in \R^p}
  \left\{
    \frac{1}{2n}\|\by - \bX\bbeta\|_2^2
    \;+\; \lambda\|\bbeta\|_1
  \right\},
  \label{eq:lasso}
\end{equation}
for a tuning parameter $\lambda > 0$. The $\ell_1$ penalty simultaneously
shrinks coefficients toward zero and sets some exactly to zero, achieving
variable selection as a by-product of estimation. Unlike best-subset
selection, the LASSO is convex and can be solved efficiently for large $p$
using coordinate descent \citep{friedman2010glmnet}.

Despite its computational appeal, the LASSO has well-documented
limitations. Variable-selection consistency requires the
\emph{irrepresentable condition} \citep{zhao2006irrepresentable}, a
stringent incoherence condition on the design matrix that is often violated
in practice. The elastic net \citep{zou2005elasticnet} and SCAD
\citep{fan2001scad} penalties address some of these shortcomings, but they
share the same fundamental limitation: they produce \emph{point estimates}
with no accompanying measure of uncertainty about which variables are truly
active. Furthermore, extending penalized methods to semi-parametric
settings requires careful treatment of the functional components.
While Group-LASSO penalties \citep{yuan2006grouplasso} can zero out entire smooth
functions, they do not support inference about \emph{localized} effects
within a non-zero function. Both limitations motivate the Bayesian
approaches developed in this dissertation.

% -----------------------------------------------------------------------------
\section{Bayesian Variable and Model Selection}
\label{sec:intro-bayes}
% -----------------------------------------------------------------------------

Bayesian statistics frames all inference as a problem of updating a prior
distribution over unknowns to a posterior distribution given observed data.
In the variable-selection context, the unknowns include not only the model
parameters $\bbeta$ but also the model \emph{structure} itself, represented
by a binary indicator vector
$\bgamma = (\gamma_1, \ldots, \gamma_p)^{\top} \in \{0,1\}^p$, where
$\gamma_j = 1$ if variable $j$ is included and $\gamma_j = 0$ otherwise.
Each configuration $\bgamma$ defines a sub-model $\mathcal{M}_{\bgamma}$,
and Bayesian variable selection assigns a probability to every such
sub-model.

\subsection{The Posterior over Model Space}
\label{subsec:posterior-model}

Given a prior $p(\bgamma)$ over model indicators and a prior
$p(\bbeta_{\bgamma} \mid \bgamma)$ over the parameters active under
$\mathcal{M}_{\bgamma}$, the joint posterior is
\begin{equation}
  p(\bgamma, \bbeta_{\bgamma} \mid \by)
  \;\propto\;
  p(\by \mid \bbeta_{\bgamma}, \bgamma)\,
  p(\bbeta_{\bgamma} \mid \bgamma)\,
  p(\bgamma).
  \label{eq:joint-posterior}
\end{equation}
Marginalizing over parameters yields the \emph{posterior model
probability} (PMP),
\begin{equation}
  p(\bgamma \mid \by)
  \;\propto\;
  p(\by \mid \bgamma)\, p(\bgamma),
  \label{eq:pmp}
\end{equation}
where $p(\by \mid \bgamma) = \int p(\by \mid \bbeta_{\bgamma},
\bgamma)\,p(\bbeta_{\bgamma}\mid\bgamma)\,\mathrm{d}\bbeta_{\bgamma}$
is the \emph{marginal likelihood} under $\mathcal{M}_{\bgamma}$. The ratio
of marginal likelihoods for two models is the \emph{Bayes factor}
\citep{kass1995bayes},
\begin{equation}
  \mathrm{BF}(\bgamma, \bgamma')
  \;=\;
  \frac{p(\by \mid \bgamma)}{p(\by \mid \bgamma')}.
  \label{eq:bayes-factor}
\end{equation}

\subsection{Priors for Variable Selection}
\label{subsec:vs-priors}

\paragraph{Spike-and-slab priors.}
The canonical Bayesian variable-selection prior is the
\emph{spike-and-slab} distribution, introduced by
\citet{mitchell1988spikeand} and developed extensively by
\citet{george1993vs} and \citet{george1997approaches}. The prior places
each coefficient under a mixture:
\begin{equation}
  \beta_j \mid \gamma_j
  \;\sim\;
  (1 - \gamma_j)\,\delta_0(\beta_j)
  \;+\;
  \gamma_j\,\N(0, \tau^2),
  \label{eq:spike-slab}
\end{equation}
where $\delta_0$ is a point mass at zero (the ``spike'') and $\N(0,\tau^2)$
is a diffuse Gaussian (the ``slab''). The indicator $\bgamma$ then follows
a prior such as $\gamma_j \overset{\mathrm{iid}}{\sim}
\mathrm{Bernoulli}(\pi)$, with $\pi$ either fixed or given a further
hyperprior. A model-complexity prior that penalizes models with many active
variables is often used to guard against overfitting \citep{scott2010bayes}.

\paragraph{Continuous shrinkage priors.}
An alternative to discrete indicator variables is to place a continuous
prior on $\bbeta$ that concentrates near zero while allowing heavy tails
for large signals. The \emph{horseshoe prior} of
\citet{carvalho2010horseshoe} achieves this through a global-local scale
mixture of Gaussians:
\begin{equation}
  \beta_j \mid \lambda_j, \tau
  \;\sim\;
  \N(0,\, \lambda_j^2 \tau^2),
  \quad
  \lambda_j \sim \mathrm{Half\text{-}Cauchy}(0,1),
  \quad
  \tau \sim \mathrm{Half\text{-}Cauchy}(0,1),
  \label{eq:horseshoe}
\end{equation}
where $\tau$ is a \emph{global} shrinkage parameter and $\lambda_j$ are
\emph{local} scales that allow individual coefficients to escape shrinkage.
The resulting marginal prior on $\beta_j$ has a pole at zero and heavy
tails, mimicking a spike-and-slab without requiring discrete model
indicators---a computational advantage that becomes especially significant
in deep neural networks (Chapter~\ref{chap:horseshoe}).

\subsection{Advantages over Frequentist Methods}
\label{subsec:bayes-advantages}

Bayesian variable selection offers several advantages over frequentist
penalized-likelihood methods.

\begin{enumerate}
  \item \textbf{Uncertainty quantification.} The \emph{posterior inclusion
  probability} (PIP), $\Prob(\gamma_j = 1 \mid \by)$, provides a
  coherent, calibrated measure of evidence for each predictor that has no
  direct frequentist analogue.

  \item \textbf{Model averaging.} Predictions and estimates can be averaged
  across models weighted by their posterior probabilities, yielding
  \emph{Bayesian model averaging} (BMA) \citep{hoeting1999bma}, which
  typically improves predictive accuracy and avoids the instability of
  single-model selection.

  \item \textbf{Automatic multiplicity control.} By penalizing large models
  through the prior on $\bgamma$, Bayesian methods naturally adjust for the
  number of comparisons being made, without requiring separate multiplicity
  corrections \citep{scott2010bayes}.

  \item \textbf{Interpretability.} The PIP is directly interpretable as a
  probability, unlike LASSO coefficient paths, which depend on the
  arbitrary scale of $\lambda$.

  \item \textbf{Coherent incorporation of prior knowledge.} Domain
  knowledge about which variables are a priori likely to be active can be
  encoded directly in $p(\bgamma)$ or in the parameter priors.
\end{enumerate}

% -----------------------------------------------------------------------------
\section{Variable Selection in Semi-Parametric Settings}
\label{sec:intro-semipar}
% -----------------------------------------------------------------------------

The methods reviewed above were developed primarily for parametric linear
models. In many contemporary applications, the relationship between
predictors and outcome is better described by a semi-parametric or
nonparametric model, such as the GAM in \eqref{eq:gam}. Extending variable
selection to these settings introduces new challenges.

\subsection{Functional Variable Selection in Generalized Additive Models}
\label{subsec:functional-vs}

In the GAM setting, variable selection reduces to testing
$H_{0j}: f_j \equiv 0$ for each $j = 1,\ldots,p$. The most common approach
represents each $f_j$ in a basis expansion,
$f_j(x) \approx \sum_{k=1}^{K} \phi_{jk}(x)\,\alpha_{jk}$,
and tests whether $\bm{\alpha}_j = (\alpha_{j1},\ldots,\alpha_{jK})^{\top}$
is zero. Basis choices include B-splines \citep{eilers1996bsplines},
natural splines, and Fourier bases.

However, standard basis expansions are not well suited to detecting
\emph{localized} effects: a smooth function $f_j$ may be non-zero only over
a compact sub-interval of its domain, yet a global basis will distribute
that local signal across many basis coefficients. As a consequence,
selection procedures based on such representations fail to achieve
asymptotic consistency for localized effects \citep{rossell2025semiparametric}.
Chapter~\ref{chap:semipar} resolves this issue by constructing an
\emph{orthogonalized} basis that eliminates overlap between basis functions,
thereby permitting each basis coefficient to capture an independent,
spatially localized contribution to $f_j$. Paired with a model-complexity
prior that penalizes unnecessary flexibility, the resulting Bayesian
procedure achieves consistent selection of localized effects even under
model misspecification.

\subsection{Variable Selection in Neural Networks}
\label{subsec:nn-vs}

Deep neural networks have emerged as a highly expressive class of function
approximators, capable of capturing complex, high-dimensional response
surfaces that resist parametric description \citep{goodfellow2016deep}.
A fully connected network with $L$ hidden layers maps an input
$\bx \in \R^p$ to a scalar output through a composition of affine
transformations and element-wise nonlinearities:
\begin{equation}
  f(\bx)
  \;=\;
  W_L \,\sigma\!\left(W_{L-1}\,\sigma\!\left(\cdots
  \sigma(W_1 \bx + \mathbf{b}_1) \cdots\right)
  + \mathbf{b}_{L-1}\right) + b_L,
  \label{eq:dnn}
\end{equation}
where $W_\ell$ and $\mathbf{b}_\ell$ are the weight matrix and bias
vector at layer $\ell$, and $\sigma$ is an activation function such as
ReLU. Input-variable selection means identifying which components of
$\bx$ are necessary for accurate prediction---equivalent to asking which
columns of $W_1$ are not identically zero.

Frequentist approaches to neural-network sparsity rely primarily on
post-hoc pruning or $\ell_1$/$\ell_0$ regularization of the weights
\citep{lecun1990obd, han2015pruning}. These methods produce sparse
networks but provide no probabilistic measure of which inputs are truly
relevant. A fully Bayesian treatment places a prior over the network
weights and updates it to a posterior given data, yielding principled
uncertainty quantification. \citet{louizos2017bayesian} showed that the
horseshoe prior of \citet{carvalho2010horseshoe} can be embedded in a
neural network via a global-local scale-mixture parameterization,
enabling automatic input-variable selection through the posterior
distribution of the local scales $\lambda_j$ at the first layer.
Chapter~\ref{chap:horseshoe} builds on and modifies this framework,
developing a tractable variational inference procedure and demonstrating
that the resulting method achieves stable, interpretable input-variable
selection across a range of nonlinear simulation settings and real
datasets.

% -----------------------------------------------------------------------------
\section{Data}
\label{sec:intro-data}
% -----------------------------------------------------------------------------

The empirical work in Chapter~\ref{chap:semipar} is illustrated on two
datasets that represent complementary challenges in semi-parametric
variable selection: a wage dataset with socioeconomic predictors, and a
high-frequency electrocorticography (ECoG) dataset from cognitive
neuroscience.

\subsection{Wage Data}
\label{subsec:data-wages}

The wage dataset is a cross-sectional sample of male workers drawn from
the \emph{Current Population Survey} \citep{murphy1990empirical}. The
response variable is the (log) hourly wage, and the predictors include
years of education, years of potential labor-market experience, and
demographic indicators. The relationship between wages and experience is
widely recognized to be nonlinear---roughly hump-shaped over the working
life---making a generalized additive model a natural choice. A key
scientific question is whether the effect of experience on wages is
\emph{globally} smooth or whether it is concentrated in specific
segments of the experience distribution, such as the early career
period. The localized variable-selection method of
Chapter~\ref{chap:semipar} is particularly well suited to answering
this question.

\subsection{Electrocorticography (ECoG) Data}
\label{subsec:data-ecog}

Electrocorticography (ECoG) records electrical activity directly from
the surface of the brain at high temporal resolution
\citep{miller2009power}. Unlike scalp electroencephalography (EEG),
ECoG electrodes are placed subdurally and thus pick up local field
potentials with far greater signal-to-noise ratio. The dataset used in
Chapter~\ref{chap:semipar} consists of broadband power measurements
recorded from a grid of electrodes placed over the cortex of a human
participant during a motor or cognitive task. Each electrode's
time-averaged power serves as a response variable, and the predictors
encode spatial or temporal features of the experimental stimuli.

The ECoG setting exemplifies the localized-effect problem: the brain's
functional organization is spatially structured, so a stimulus may
activate only a subset of electrodes, and the effect may be restricted
to a particular frequency band or time window. Detecting such localized
activations while controlling for the large number of electrodes tested
simultaneously is precisely the inferential challenge addressed in
Chapter~\ref{chap:semipar}.

% -----------------------------------------------------------------------------
\section{Dissertation Outline}
\label{sec:intro-outline}
% -----------------------------------------------------------------------------

The remainder of this dissertation is organized as follows.

\textbf{Chapter~\ref{chap:semipar}} presents the work of
\citet{rossell2025semiparametric} on semiparametric local variable
selection under misspecification. We introduce a model-complexity prior
for generalized additive models, construct an orthogonalized basis that
eliminates basis-function overlap, and prove that the resulting Bayesian
selection procedure is asymptotically consistent for localized
functional effects even when the working model is misspecified.
Simulation studies illustrate the method's advantages over standard
basis-expansion competitors, and the wage and ECoG datasets are
analyzed in detail.

\textbf{Chapter~\ref{chap:horseshoe}} develops a Bayesian neural
network for input-variable selection based on a modified horseshoe
prior. We describe the model architecture, the variational inference
algorithm used for posterior approximation, and theoretical properties
of the induced shrinkage. The method is evaluated on simulated datasets
with known nonlinear structure and on real-data benchmarks.

\textbf{The Appendix} collects supplementary material including
additional proofs, simulation tables, and implementation details for
both chapters.

%%% Local Variables: ***
%%% mode: latex ***
%%% TeX-master: "thesis.tex" ***
%%% End: ***